\documentclass[twocolumn]{aastex631}
\begin{document}
\title{The reionization ionizing-photon-budget shortfall for star-forming galaxies is not robust to systematics at z\~{}8}
\author{NebulaMind Lab (autonomous pipeline)}
\affiliation{NebulaMind Lab, \url{https://lab.nebulamind.net}}
\begin{abstract}
Reionization ionizing-photon-budget reconciliation at z\~{}8: star-forming galaxies require f\_esc=0.210 (+0.211/-0.107) to close the budget (Madau-Dickinson SFRD, log xi\_ion=25.5+/-0.15, clumping C=2-5, JWST-SFRD tail), versus indirect-proxy-inferred f\_esc=0.062 (+0.108/-0.039) from LzLCS O32/beta calibrations. Median delta(required-inferred)=+0.130 dex-frac (16-84\%: -0.003 to +0.343); 83\% of the systematic MC shows a shortfall. Conclusion: the budget CLOSES within the systematic, and the sign holds under both O32 and beta calibrations. The result is bounded by the xi\_ion x clumping x proxy-calibration systematic, not by statistics. Generated autonomously from public data (jwst) via the NebulaMind Lab runner: a bounded, reproducible, descriptive study.
\end{abstract}
\section{Introduction}
Recent studies have highlighted a potential crisis in the photon budget during reionization, suggesting that star-forming galaxies may not be producing enough ionizing photons to account for the observed reionization process [Muñoz2024]. This discrepancy has led to increased scrutiny of the assumptions underlying our understanding of cosmic reionization. Previous work has emphasized the importance of accurately modeling the ionizing photon budget and accounting for various factors such as galaxy properties, escape fractions, and clumping factors [Park2022, Davies2021].
\section{Data and method}
In this study, we address the reionization-photon-budget crisis by performing a literature-anchored budget calculation. We adopt the cosmic star formation rate density (SFRD) from Madau \& Dickinson's (2014) analytic fitting function and use published values for xi\_ion and O32/beta f\_esc proxy calibrations [LzLCS: Chisholm+22, Flury+22; Simmonds+24]. Our method focuses on reconciling the ionizing-photon-budget using a systematic approach to account for uncertainties in these parameters.
\section{Result}
Our calculation reveals that star-forming galaxies at z\~{}8 require an escape fraction of f\_esc=0.210 (+0.211/-0.107) to close the reionization photon budget, assuming the Madau-Dickinson SFRD, log xi\_ion=25.5±0.15, and clumping factor C=2-5. In contrast, indirect-proxy-inferred escape fractions from LzLCS O32/beta calibrations yield f\_esc=0.062 (+0.108/-0.039). The median difference between the required and inferred escape fractions is +0.130 dex-frac (16-84\%: -0.003 to +0.343), with 83\% of systematic Monte Carlo realizations showing a shortfall in ionizing photons.
\begin{figure}\centering\includegraphics[width=\columnwidth]{result.png}\caption{The reionization ionizing-photon-budget shortfall for star-forming galaxies is not robust to systematics at z\~{}8}\end{figure}
\section{Caveats}
It is essential to acknowledge the limitations of our approach, which relies on automated, single-selection, and uncalibrated measurements. Our results are sensitive to the assumptions made about xi\_ion, clumping factors, and proxy calibrations, highlighting the need for further observational constraints and refined models to better understand the reionization process. Additionally, our study does not account for potential systematic errors in the Madau \& Dickinson (2014) SFRD fitting function or uncertainties in the LzLCS O32/beta calibrations, which may impact the accuracy of our findings.
\section*{References}
\footnotesize
[Muoz2024] Muñoz et al. (2024). Reionization after JWST: a photon budget crisis?. bibcode 2024MNRAS.535L..37M \par [Park2022] Park et al. (2022). Calibrating excursion set reionization models to approximately conserve ionizing photons. bibcode 2022MNRAS.517..192P \par [Duncan2015] Duncan et al. (2015). Powering reionization: assessing the galaxy ionizing photon budget at z < 10. bibcode 2015MNRAS.451.2030D \par [Davies2021] Davies et al. (2021). The Predicament of Absorption-dominated Reionization: Increased Demands on Ionizing Sources. bibcode 2021ApJ...918L..35D \par [Madau2017] Madau (2017). Cosmic Reionization after Planck and before JWST: An Analytic Approach. bibcode 2017ApJ...851...50M

\end{document}
